% Options for packages loaded elsewhere
% Options for packages loaded elsewhere
\PassOptionsToPackage{unicode}{hyperref}
\PassOptionsToPackage{hyphens}{url}
\PassOptionsToPackage{dvipsnames,svgnames,x11names}{xcolor}
%
\documentclass[
  11pt]{article}
\usepackage{xcolor}
\usepackage[margin=1in]{geometry}
\usepackage{amsmath,amssymb}
\setcounter{secnumdepth}{5}
\usepackage{iftex}
\ifPDFTeX
  \usepackage[T1]{fontenc}
  \usepackage[utf8]{inputenc}
  \usepackage{textcomp} % provide euro and other symbols
\else % if luatex or xetex
  \usepackage{unicode-math} % this also loads fontspec
  \defaultfontfeatures{Scale=MatchLowercase}
  \defaultfontfeatures[\rmfamily]{Ligatures=TeX,Scale=1}
\fi
\usepackage{lmodern}
\ifPDFTeX\else
  % xetex/luatex font selection
\fi
% Use upquote if available, for straight quotes in verbatim environments
\IfFileExists{upquote.sty}{\usepackage{upquote}}{}
\IfFileExists{microtype.sty}{% use microtype if available
  \usepackage[]{microtype}
  \UseMicrotypeSet[protrusion]{basicmath} % disable protrusion for tt fonts
}{}
\makeatletter
\@ifundefined{KOMAClassName}{% if non-KOMA class
  \IfFileExists{parskip.sty}{%
    \usepackage{parskip}
  }{% else
    \setlength{\parindent}{0pt}
    \setlength{\parskip}{6pt plus 2pt minus 1pt}}
}{% if KOMA class
  \KOMAoptions{parskip=half}}
\makeatother
% Make \paragraph and \subparagraph free-standing
\makeatletter
\ifx\paragraph\undefined\else
  \let\oldparagraph\paragraph
  \renewcommand{\paragraph}{
    \@ifstar
      \xxxParagraphStar
      \xxxParagraphNoStar
  }
  \newcommand{\xxxParagraphStar}[1]{\oldparagraph*{#1}\mbox{}}
  \newcommand{\xxxParagraphNoStar}[1]{\oldparagraph{#1}\mbox{}}
\fi
\ifx\subparagraph\undefined\else
  \let\oldsubparagraph\subparagraph
  \renewcommand{\subparagraph}{
    \@ifstar
      \xxxSubParagraphStar
      \xxxSubParagraphNoStar
  }
  \newcommand{\xxxSubParagraphStar}[1]{\oldsubparagraph*{#1}\mbox{}}
  \newcommand{\xxxSubParagraphNoStar}[1]{\oldsubparagraph{#1}\mbox{}}
\fi
\makeatother


\usepackage{longtable,booktabs,array}
\usepackage{calc} % for calculating minipage widths
% Correct order of tables after \paragraph or \subparagraph
\usepackage{etoolbox}
\makeatletter
\patchcmd\longtable{\par}{\if@noskipsec\mbox{}\fi\par}{}{}
\makeatother
% Allow footnotes in longtable head/foot
\IfFileExists{footnotehyper.sty}{\usepackage{footnotehyper}}{\usepackage{footnote}}
\makesavenoteenv{longtable}
\usepackage{graphicx}
\makeatletter
\newsavebox\pandoc@box
\newcommand*\pandocbounded[1]{% scales image to fit in text height/width
  \sbox\pandoc@box{#1}%
  \Gscale@div\@tempa{\textheight}{\dimexpr\ht\pandoc@box+\dp\pandoc@box\relax}%
  \Gscale@div\@tempb{\linewidth}{\wd\pandoc@box}%
  \ifdim\@tempb\p@<\@tempa\p@\let\@tempa\@tempb\fi% select the smaller of both
  \ifdim\@tempa\p@<\p@\scalebox{\@tempa}{\usebox\pandoc@box}%
  \else\usebox{\pandoc@box}%
  \fi%
}
% Set default figure placement to htbp
\def\fps@figure{htbp}
\makeatother


% definitions for citeproc citations
\NewDocumentCommand\citeproctext{}{}
\NewDocumentCommand\citeproc{mm}{%
  \begingroup\def\citeproctext{#2}\cite{#1}\endgroup}
\makeatletter
 % allow citations to break across lines
 \let\@cite@ofmt\@firstofone
 % avoid brackets around text for \cite:
 \def\@biblabel#1{}
 \def\@cite#1#2{{#1\if@tempswa , #2\fi}}
\makeatother
\newlength{\cslhangindent}
\setlength{\cslhangindent}{1.5em}
\newlength{\csllabelwidth}
\setlength{\csllabelwidth}{3em}
\newenvironment{CSLReferences}[2] % #1 hanging-indent, #2 entry-spacing
 {\begin{list}{}{%
  \setlength{\itemindent}{0pt}
  \setlength{\leftmargin}{0pt}
  \setlength{\parsep}{0pt}
  % turn on hanging indent if param 1 is 1
  \ifodd #1
   \setlength{\leftmargin}{\cslhangindent}
   \setlength{\itemindent}{-1\cslhangindent}
  \fi
  % set entry spacing
  \setlength{\itemsep}{#2\baselineskip}}}
 {\end{list}}
\usepackage{calc}
\newcommand{\CSLBlock}[1]{\hfill\break\parbox[t]{\linewidth}{\strut\ignorespaces#1\strut}}
\newcommand{\CSLLeftMargin}[1]{\parbox[t]{\csllabelwidth}{\strut#1\strut}}
\newcommand{\CSLRightInline}[1]{\parbox[t]{\linewidth - \csllabelwidth}{\strut#1\strut}}
\newcommand{\CSLIndent}[1]{\hspace{\cslhangindent}#1}



\setlength{\emergencystretch}{3em} % prevent overfull lines

\providecommand{\tightlist}{%
  \setlength{\itemsep}{0pt}\setlength{\parskip}{0pt}}



 


\makeatletter
\@ifpackageloaded{caption}{}{\usepackage{caption}}
\AtBeginDocument{%
\ifdefined\contentsname
  \renewcommand*\contentsname{Table of contents}
\else
  \newcommand\contentsname{Table of contents}
\fi
\ifdefined\listfigurename
  \renewcommand*\listfigurename{List of Figures}
\else
  \newcommand\listfigurename{List of Figures}
\fi
\ifdefined\listtablename
  \renewcommand*\listtablename{List of Tables}
\else
  \newcommand\listtablename{List of Tables}
\fi
\ifdefined\figurename
  \renewcommand*\figurename{Figure}
\else
  \newcommand\figurename{Figure}
\fi
\ifdefined\tablename
  \renewcommand*\tablename{Table}
\else
  \newcommand\tablename{Table}
\fi
}
\@ifpackageloaded{float}{}{\usepackage{float}}
\floatstyle{ruled}
\@ifundefined{c@chapter}{\newfloat{codelisting}{h}{lop}}{\newfloat{codelisting}{h}{lop}[chapter]}
\floatname{codelisting}{Listing}
\newcommand*\listoflistings{\listof{codelisting}{List of Listings}}
\makeatother
\makeatletter
\makeatother
\makeatletter
\@ifpackageloaded{caption}{}{\usepackage{caption}}
\@ifpackageloaded{subcaption}{}{\usepackage{subcaption}}
\makeatother
\usepackage{bookmark}
\IfFileExists{xurl.sty}{\usepackage{xurl}}{} % add URL line breaks if available
\urlstyle{same}
\hypersetup{
  pdftitle={Bell Violations as Three-Party Quadratic Conservation: A Classical Three-Phase Realization},
  pdfauthor={David Nobles},
  pdfkeywords={Bell inequalities, Tsirelson bound, Clarke
transform, prequantum classical models, quantum foundations, three-phase
power systems},
  colorlinks=true,
  linkcolor={blue},
  filecolor={Maroon},
  citecolor={Blue},
  urlcolor={Blue},
  pdfcreator={LaTeX via pandoc}}


\title{Bell Violations as Three-Party Quadratic Conservation: A
Classical Three-Phase Realization}
\author{David Nobles}
\date{2026-05-04}
\begin{document}
\maketitle
\begin{abstract}
The hierarchy of Bell-inequality violations arises from the algebraic
order of the conservation law governing the correlated source. Linear
conservation (momentum-type) produces correlations bounded by
\(S = \sqrt{2}\), well within the classical Bell--CHSH limit. Quadratic
conservation (power-, intensity-, or \(|\psi|^2\)-type) produces
correlations that saturate the Tsirelson bound at \(S = 2\sqrt{2}\). The
factor-of-two amplification in angular frequency ---
\(\cos(2\Delta\phi)\) versus \(\cos(\Delta\phi)\) --- is the
double-angle identity made physical. We argue that the standard
two-party CHSH framework constitutes an incomplete accounting of an
intrinsically three-party conservation law, with the third party
identified as the zero-sequence mode of the Clarke decomposition: the
energy carried by the field between measurement stations, structurally
analogous to vacuum and edge-mode contributions in algebraic quantum
field theory and to the quantum potential in Bohmian mechanics. We
construct a classical three-phase analog system --- a passive linear
front-end of four delta-LC tanks with a shared resonant core, a
shot-trigger Born-rule detector, and a zero-sequence closure latch ---
and show by reduced-model simulation and a SPICE-driven realization that
it reproduces Tsirelson-saturating CHSH (\(S \approx 2.828\), error
\(0.043\); winner-law RMS error \(0.013\); decisive rate \(99.5\%\);
energy conservation \(100\%\)). We position this construction relative
to prequantum classical models (Khrennikov's PCSFT, Spekkens,
't\textasciitilde Hooft, Bohmian mechanics) and clarify scope: the
system requires a non-local conservation constraint and therefore does
\emph{not} refute Bell's theorem. Its contribution is to make the locus
of nonlocality concrete --- the zero-sequence energy flow --- and to
give the apparent nonlocality of quantum correlations a structural
interpretation in terms of three-party quadratic conservation
accounting.
\end{abstract}


\section{Introduction}\label{sec-intro}

Bell's theorem (Bell 1964) establishes that the correlations between
spacelike-separated measurements predicted by quantum mechanics cannot
be reproduced by any local hidden variable (LHV) model. The
Clauser--Horne--Shimony--Holt (CHSH) inequality (Clauser et al. 1969)
gives the standard quantitative test: for any LHV model, the combination
\[
S \;=\; E(a,b) - E(a,b') + E(a',b) + E(a',b')
\] satisfies \(|S| \leq 2\), where \(E(a,b)\) is the expectation of the
product of binary \(\pm 1\) outcomes at analyzer settings \(a\) and
\(b\). Quantum mechanics predicts \(|S| \leq 2\sqrt{2}\), the Tsirelson
bound (Tsirelson 1980), saturated by singlet-state measurements at
appropriately chosen angles. Loophole-free experiments (Hensen et al.
2015; Giustina et al. 2015; Shalm et al. 2015) have confirmed the
violation under conditions that close the locality and detection
loopholes simultaneously.

The standard interpretation frames the violation as evidence that nature
is fundamentally nonlocal, or that realism --- the assignment of
pre-existing values to unmeasured observables --- must be abandoned.
This paper offers a complementary interpretation grounded in the algebra
of conservation laws and a constructive demonstration in classical
three-phase circuits. The two are presented in lockstep because each,
taken alone, invites the wrong reading.

\subsection{The interpretive thesis}\label{the-interpretive-thesis}

We claim that the form of Bell-correlation functions is determined by
the algebraic order of the conservation law governing the source, and
that the standard CHSH framework systematically omits a third
participant in what is intrinsically a tripartite conservation
constraint.

When the conserved quantity at the source is \emph{linear} in the field
amplitude --- as in two-body decays where momentum
\(\mathbf{p}_A + \mathbf{p}_B = 0\) --- the joint correlation has the
form \(E(a,b) \propto \cos(\Delta\phi)\), and the CHSH combination
satisfies \(|S| \leq \sqrt{2}\), well within the classical limit. When
the conserved quantity is \emph{quadratic} --- power, intensity,
\(|\psi|^2\), or any \(\sum_k v_k^2\) form --- the correlation becomes
\(E(a,b) \propto \cos(2\Delta\phi)\), and CHSH attains
\(|S| = 2\sqrt{2}\), saturating Tsirelson. The factor of two inside the
cosine argument is not an accident of the quantum formalism. It is the
angle-doubling inherent in squaring: \(\cos^2(x) = (1 + \cos(2x))/2\).
The passage from amplitude to intensity, from \(\psi\) to \(|\psi|^2\),
doubles the angular frequency of the correlation pattern, and that is
the algebraic engine that pushes \(S\) past the Bell bound.

This re-reading of the cosine doubling has a structural consequence. A
quadratic conservation law on a vector field with rotational symmetry is
naturally a \emph{three-party} law, even when only two parties are
measured. The Clarke transform (Clarke 1943) from three-phase power
systems decomposes any three-component balanced signal
\((v_a, v_b, v_c)\) into orthogonal components
\((v_\alpha, v_\beta, v_0)\), and Parseval's identity gives \[
v_a^2 + v_b^2 + v_c^2 \;=\; v_\alpha^2 + v_\beta^2 + v_0^2 \;=\; \mathrm{const}.
\] The two analyzers in a CHSH experiment project onto two of the three
Clarke modes; the third --- the zero-sequence mode \(v_0\) --- is energy
that flows neither to Alice nor to Bob, but \emph{between} them. We
identify this zero-sequence channel as the third party in the
conservation law.

Concretely: if \(P_A(\theta)\) and \(P_B(\theta)\) are the powers
captured by analyzers at settings \(\phi_A\) and \(\phi_B\) on a
balanced positive--negative-sequence source with hidden emission phase
\(\theta\), then a direct computation (Section~\ref{sec-hierarchy})
gives \[
P_A(\theta) + P_B(\theta) \;=\; 1 + \cos(\Delta\phi)\cos(2\theta - \phi_A - \phi_B),
\] which is \emph{not} constant unless \(\Delta\phi \in \{0, \pi/2\}\).
The remaining energy resides in the zero-sequence mode, \[
P_0(\theta) \;=\; P_\text{total} - P_A(\theta) - P_B(\theta),
\] which oscillates in \(\theta\) at exactly the angular frequency that
makes \(P_+ + P_- + P_0\) constant. The CHSH violation, on this reading,
measures the depth of the zero-sequence reservoir at the
analyzer-difference angle \(\Delta\phi = \pi/4\). Aligned analyzers
(\(\Delta\phi = 0\)) drain \(P_0\) to zero; orthogonal analyzers
(\(\Delta\phi = \pi/2\)) push all correlation energy into \(P_0\);
intermediate angles are precisely where the Bell bound is exceeded.

The structural analogy extends naturally to algebraic quantum field
theory, where the entanglement between two spatial regions is carried by
the field modes at their boundary --- the ``edge modes'' or ``zero
modes'' that belong to neither region's local algebra (Haag 1996; Witten
2018). In Bohmian mechanics, the same structural role is played by the
quantum potential, the nonlocal term in the guidance equation that
encodes the global state's constraint on local dynamics (Bohm 1952;
Holland 1993). The zero-sequence mode in the Clarke decomposition is, on
the present reading, the explicit physical reservoir corresponding to
those abstractions in a setting where every degree of freedom is
classical and accessible to direct measurement.

\subsection{The constructive thesis}\label{the-constructive-thesis}

The interpretive thesis can be stated abstractly. To make it concrete,
and to test it where it can be tested, we construct a classical
three-phase analog system that implements the full pipeline from source
state preparation through analyzer projection through detection and
post-click closure, and demonstrate by simulation that it reproduces
Tsirelson-saturating CHSH.

The system has four layers. (i) A reduced four-mode shared state
\(\Psi_0 \in \mathbb{C}^4\) in the joint basis
\(\{|{++}\rangle, |{+-}\rangle, |{-+}\rangle, |{--}\rangle\}\), prepared
as the singlet projection. (ii) Local analyzer rotations
\(R_A(a) \otimes R_B(b)\) producing branch amplitudes
\(c_{xy}(a,b) = u_{xy}^\dagger \Psi_0\) and branch energies
\(w_{xy}(a,b) = |c_{xy}(a,b)|^2\) at the linear energy-fraction level
(Theorem~\textbf{?@thm-branch-weights}). (iii) A noisy energy-threshold
first-arrival detector that converts continuous branch energies into
discrete \(\pm 1\) click outcomes with Born-rule statistics under
explicit and disclosed conditions (Theorem~\textbf{?@thm-born}). (iv) A
zero-sequence closure latch that suppresses losing branches and routes
residual energy to the winning drain via the common-mode rail.

Each layer is implemented in code. The reduced 4-mode model lives in
\texttt{src/shared\_4tank\_core.py} and \texttt{src/joint\_readout.py}.
The detector family is searched and characterized in
\texttt{detector\_search/}. The full integration --- front-end,
detector, latch --- is exercised in \texttt{detector\_integration/}. The
SPICE-driven physical realization, in which the four delta-LC tanks
become explicit inductor--capacitor networks coupled by a shared
resonant core and the analyzer rotations are implemented as physical
coupling matrices, lives in \texttt{physical\_front\_end\_candidate/}.

The numerical results from the full chain are the headline of the
constructive contribution. The preferred LC-coupled physical chain
achieves CHSH \(S \approx 2.828\) (absolute error \(0.043\) against the
target \(2\sqrt{2}\)), winner-law RMS error \(0.013\) (max \(0.039\)),
correlator RMS error \(0.024\), decisive (one-click) rate \(99.5\%\),
winner-drain dominance \(99.97\%\), and energy conservation at \(100\%\)
within the numerical tolerance of the SPICE engine. The decisive rate is
well above the fair-sampling threshold, so the result does not depend on
the detection loophole (Larsson 2014). Every quantitative claim in the
paper is anchored to a Make target and a reproducible artifact;
Appendix~\textbf{?@sec-repro} gives the full mapping.

\subsection{What this paper does not
claim}\label{what-this-paper-does-not-claim}

A construction that reproduces Tsirelson-saturating CHSH in a simulated
classical circuit invites a misreading: that we have refuted Bell's
theorem by exhibiting a local hidden variable model with
\(|S| = 2\sqrt{2}\). That misreading is wrong, and it is important to
disclose precisely why before going further.

Bell's theorem is a statement about LHV models satisfying a specific
factorization condition: the joint outcome distribution must factor as
\(P(A,B \mid a, b, \lambda) = P(A \mid a, \lambda)\, P(B \mid b, \lambda)\)
when conditioned on the hidden variable \(\lambda\). Our construction
does not satisfy this factorization. The zero-sequence mode \(P_0\) is a
global degree of freedom that couples Alice's and Bob's measurement
statistics through the conserved quadratic form
\(P_+ + P_- + P_0 = \mathrm{const}\). When Alice's analyzer captures
more or less power, the zero-sequence reservoir absorbs or releases the
difference, and Bob's marginal is affected through that coupling --- not
through a signaling channel (the local marginals remain uniform;
Section~\ref{sec-two-party}) but through the joint constraint imposed by
the conservation law.

In Bell's terminology, the zero-sequence mode is a non-local ``common
cause'' that mediates the correlations. In the field-theoretic
terminology of (Witten 2018), it is the boundary mode shared between the
two regions. In Bohmian terminology, it is the quantum-potential
analogue. Our system reproduces the cos-squared correlation structure
--- and therefore Tsirelson saturation --- precisely because the
conservation law is enforced event-by-event across the entire analog
network. That enforcement is what no factorized local model can deliver,
and it is also what makes our construction \emph{not} an LHV model in
the technical sense Bell defined.

The paper's contribution is therefore not to refute Bell. It is to make
the locus of nonlocality concrete and physical. The ``spookiness'' of
Bell violations, on the present reading, is not an irreducible feature
of distant measurement events. It is the inevitable consequence of
incomplete accounting: the two-party CHSH framework projects out the
third channel through which the conservation law does its work. When
that channel is restored to the bookkeeping, the correlations become as
natural as momentum conservation in a two-body decay, but enriched by
exactly the factor of two that distinguishes amplitude from intensity.

\subsection{Relation to prior
frameworks}\label{relation-to-prior-frameworks}

The view developed here intersects with several established programs in
foundations.

\emph{Algebraic quantum field theory} identifies entanglement between
spatial regions as a property of boundary modes shared between them; see
(Haag 1996) for the foundational treatment and (Witten 2018) for a
recent overview. The zero-sequence mode plays exactly the structural
role those treatments assign to the field at the boundary: a degree of
freedom that is not part of either region's local observable algebra but
that is essential to the joint state's correlation structure.

\emph{Bohmian mechanics} and pilot-wave theory (Bohm 1952; Holland 1993)
supplement the wavefunction with an explicit configuration variable
guided by a nonlocal potential. The quantum potential plays the same
coupling role our zero-sequence mode plays: it is the term that enforces
the global constraint on local dynamics. Our framework is consistent
with --- though not dependent on --- this reading.

\emph{Khrennikov's prequantum classical statistical field theory}
(PCSFT) (Khrennikov 2014) models quantum observables as quadratic
functionals of classical random fields. Our construction shares PCSFT's
commitment to classical fields and quadratic functionals as the source
of quantum statistics, but differs in two respects. First, we use a
finite-dimensional lumped-element power network rather than an
infinite-dimensional functional integral, which makes every quantity
directly measurable in a laboratory. Second, we make the role of the
zero-sequence mode explicit: in PCSFT, the analog of \(P_0\) is buried
in the regularization of the field state, whereas in our construction it
is a direct physical degree of freedom on a third wire.

\emph{Spekkens's toy theory} (Spekkens 2007) reproduces a substantial
fragment of quantum behavior under an epistemic restriction in a
discrete model. The toy theory is a different project --- it shows what
\emph{can} be reproduced from a knowledge-restriction principle alone
--- and it does not produce a continuous \(\cos^2\) structure or
saturate Tsirelson. The two programs are complementary: Spekkens
isolates which features come from epistemic restriction; our
construction isolates which features come from quadratic conservation.

\emph{'t Hooft's cellular automaton interpretation} (Hooft 2016)
proposes a fundamental, deterministic, sub-quantum theory. Our framework
is structural rather than fundamental: we are not asserting that quantum
mechanics emerges from three-phase power systems, only that the algebra
of three-phase quadratic conservation provides the natural mathematical
habitat for the structure of Bell correlations.

\emph{Detection-loophole and time-dependent LHV literatures} (Larsson
2014; Hess and Philipp 2001) address whether classical models
supplemented with realistic detection inefficiencies or time-correlated
hidden variables can mimic Bell violations. The constructive part of our
work is independent of these arguments: the SPICE-driven realization
achieves a \(99.5\%\) decisive rate, well above any reasonable
fair-sampling threshold, and the cos-squared correlator emerges from the
static algebraic structure of the conservation law, not from
time-dependent randomness.

A more detailed engagement with each of these literatures appears in
Section~\textbf{?@sec-prior-art}.

\subsection{Roadmap}\label{roadmap}

Section~\ref{sec-hierarchy} establishes the conservation hierarchy:
linear conservation gives \(S = \sqrt{2}\), quadratic conservation gives
\(S = 2\sqrt{2}\), and the deterministic-threshold model that produces
the Bell bound \(S = 2\) sits between them in a precise algebraic sense.
Section~\ref{sec-clarke} introduces the Clarke transform as the natural
measurement apparatus for quadratic-conservation sources and establishes
the geometric correspondence between the Clarke power surface and the
Born-rule probability distribution. Section~\ref{sec-two-party} derives
the two-party CHSH correlation from a balanced three-phase source and
confronts the threshold and binarization problem.
Section~\ref{sec-third-party} introduces the zero-sequence mode as the
third party in the conservation law, computes its profile across a CHSH
angle sweep, and develops the analogies to AQFT edge modes and the
Bohmian quantum potential. Section~\textbf{?@sec-reduced-model}
specifies the reduced four-mode shared-state model and proves the
squared-projection law at the linear energy-fraction level.
Section~\textbf{?@sec-born} derives the Born rule from a noisy
energy-threshold first-arrival detector under explicit conditions, and
discusses the family search that selects shot-trigger nucleation.
Section~\textbf{?@sec-latch} specifies the zero-sequence closure latch
and characterizes its post-click behavior. Section~\textbf{?@sec-spice}
presents the SPICE-driven physical realization and its numerical
results. Section~\textbf{?@sec-prior-art} positions the construction
against prior frameworks in detail. Section~\textbf{?@sec-discussion}
discusses predictions, limitations, and open questions, and
Section~\textbf{?@sec-conclusion} concludes.

\section{The Conservation Hierarchy}\label{sec-hierarchy}

We begin by establishing the central algebraic claim: the form of the
joint correlation function \(E(a,b)\) is determined by the order of the
conservation law on the source. Linear conservation produces
\(E \propto \cos(\Delta\phi)\) and CHSH bounded by \(\sqrt{2}\).
Quadratic conservation produces \(E \propto \cos(2\Delta\phi)\) and CHSH
bounded by \(2\sqrt{2}\). The angle doubling between the two cases is
the double-angle identity \(\cos^2(x) = (1 + \cos(2x))/2\) promoted from
arithmetic to physics. The deterministic-threshold model that produces
the Bell bound \(S = 2\) exactly --- and no violation --- sits between
the two cases in a way that we make algebraically precise.

\subsection{General setup}\label{general-setup}

Consider a source that emits two correlated subsystems \(A\) and \(B\)
propagating to spatially separated measurement stations. Each station
applies a projection at an analyzer angle (\(\phi_A\) for Alice,
\(\phi_B\) for Bob) chosen independently. The joint outcome is
correlated through a hidden source variable \(\theta\), taken throughout
to be the emission phase, distributed uniformly on \([0, 2\pi)\). The
correlation function is \[
E(\phi_A, \phi_B) \;=\; \int_0^{2\pi} m_A(\theta, \phi_A)\, m_B(\theta, \phi_B)\, \frac{d\theta}{2\pi},
\] where \(m_A\) and \(m_B\) are the (continuous-valued) outputs of the
two analyzers. The form of \(m_A\) and \(m_B\) is fixed by the physics
of the source and the analyzer; in particular, by the structure of the
conservation law on the field carried from source to analyzer.

We examine two cases.

\subsection{Case I: linear conservation
(momentum-type)}\label{case-i-linear-conservation-momentum-type}

A particle at rest decays into two fragments with momenta satisfying
\(\mathbf{p}_A + \mathbf{p}_B = 0\). Each fragment carries a definite
momentum vector determined at the decay event and labelled by an
unobserved emission angle \(\theta\). An analyzer at angle \(\phi\)
projects the momentum onto the measurement axis: \[
m_A(\theta, \phi_A) = \cos(\theta - \phi_A), \qquad m_B(\theta, \phi_B) = -\cos(\theta - \phi_B),
\] where the minus sign on \(m_B\) enforces the anti-correlation imposed
by \(\mathbf{p}_A = -\mathbf{p}_B\). The correlation is then \[
E_\text{lin}(\phi_A, \phi_B) \;=\; -\int_0^{2\pi} \cos(\theta - \phi_A)\cos(\theta - \phi_B)\, \frac{d\theta}{2\pi}.
\] Applying the product-to-sum identity
\(\cos x \cos y = \tfrac{1}{2}[\cos(x-y) + \cos(x+y)]\) and using the
fact that
\(\int_0^{2\pi} \cos(2\theta - \phi_A - \phi_B)\, d\theta = 0\), we
obtain \[
E_\text{lin}(\phi_A, \phi_B) \;=\; -\tfrac{1}{2}\cos(\Delta\phi), \qquad \Delta\phi := \phi_A - \phi_B.
\]

For the CHSH evaluation, the optimal analyzer settings for a
single-frequency cosine correlator \(-k\cos(\omega\Delta\phi)\) are at
angle differences
\(\Delta\phi \in \{\pi/(4\omega), -\pi/(4\omega), -\pi/(4\omega), -\pi/(4\omega)\}\),
which yields \(|S|_{\max} = 2\sqrt{2}\, k\). For the linear correlator
\(\omega = 1\), this gives the standard choice \(\phi_A = 0\),
\(\phi_{A'} = \pi/2\), \(\phi_B = \pi/4\), \(\phi_{B'} = 3\pi/4\).
Direct substitution gives \[
S_\text{lin} \;=\; \big| E_\text{lin}(\phi_A, \phi_B) - E_\text{lin}(\phi_A, \phi_{B'}) + E_\text{lin}(\phi_{A'}, \phi_B) + E_\text{lin}(\phi_{A'}, \phi_{B'}) \big| \;=\; \sqrt{2} \;\approx\; 1.414.
\] This is well below the Bell bound \(|S| \leq 2\). Linear conservation
produces no Bell violation, and indeed sits below the classical limit by
a factor of \(\sqrt{2}\).

\subsection{Case II: quadratic conservation
(power-type)}\label{case-ii-quadratic-conservation-power-type}

Now consider a source whose conserved quantity is quadratic in the
fundamental field variable. In an electrical realization, the conserved
quantity is power \(P = V^2/R\) rather than voltage \(V\). In a quantum
realization, it is probability density \(|\psi|^2\) rather than
amplitude \(\psi\). In an electromagnetic realization, it is intensity
\(|\mathbf{E}|^2\) rather than the field \(\mathbf{E}\). In each case,
the source emits two counter-rotating components --- positive and
negative sequence in the electrical case, two helicities in the EM case,
two amplitude branches in the quantum case --- and the analyzer extracts
the power on its projection axis.

For a positive-sequence signal of amplitude \(V\) measured by an
analyzer at angle \(\phi\), the captured power is the squared
projection, \[
P_+(\theta, \phi) \;=\; V^2 \cos^2(\theta - \phi),
\] and the complementary negative-sequence power is \[
P_-(\theta, \phi) \;=\; V^2 \sin^2(\theta - \phi).
\] Normalize and centre the analyzer outputs as
\(A(\theta, \phi_A) = 2 P_+(\theta, \phi_A)/V^2 - 1\) and
\(B(\theta, \phi_B) = 1 - 2 P_+(\theta, \phi_B)/V^2\) to put them on the
\(\pm 1\) scale used in the CHSH evaluation. Then
\(A = \cos(2(\theta - \phi_A))\) and \(B = -\cos(2(\theta - \phi_B))\),
and \[
E_\text{quad}(\phi_A, \phi_B) \;=\; -\int_0^{2\pi} \cos(2(\theta - \phi_A)) \cos(2(\theta - \phi_B))\, \frac{d\theta}{2\pi} \;=\; -\tfrac{1}{2}\cos(2\Delta\phi).
\] The factor of two inside the cosine is the algebraic engine of the
violation. It arises directly from the double-angle identity \[
\cos^2(x) \;=\; \frac{1 + \cos(2x)}{2}.
\] The squaring operation --- the passage from amplitude to power, from
\(\psi\) to \(|\psi|^2\) --- doubles the angular frequency at which the
correlator oscillates against analyzer-difference angle. Whatever
happens at \(\Delta\phi\) in the linear case happens at \(2\Delta\phi\)
in the quadratic case.

For the CHSH evaluation, the doubled angular frequency \(\omega = 2\)
halves the optimal analyzer-angle spacing to \(\phi_A = 0\),
\(\phi_{A'} = \pi/4\), \(\phi_B = \pi/8\), \(\phi_{B'} = 3\pi/8\). With
strict event-by-event complementarity (Section~\ref{sec-two-party}), the
correlator amplitude is unity and direct substitution gives \[
S_\text{quad} \;=\; 2\sqrt{2} \;\approx\; 2.828.
\] This saturates the Tsirelson bound. The violation \(S > 2\) is the
\emph{direct algebraic consequence} of the conservation quantity being
quadratic rather than linear in the field; correspondingly, the optimal
CHSH analyzer-angle scale halves --- the angle scale at which violation
peaks shrinks by exactly the factor of two by which the correlator's
angular frequency doubles. The same analyzer geometry, the same
hidden-variable distribution, the same factorization of the source into
two counter-rotating components --- only the algebraic order of the
conserved quantity differs --- and the CHSH value moves from
\(\sqrt{2}\) to \(2\sqrt{2}\).

The pre-normalization expression
\(E_\text{quad} = -\tfrac{1}{2}\cos(2\Delta\phi)\) has amplitude
\(1/2\), not \(1\), and produces \(S = \sqrt{2}\) on its own; the
doubling to \(2\sqrt{2}\) comes from the strict event-by-event
complementarity \(\cos^2(\theta - \phi) + \sin^2(\theta - \phi) = 1\)
that fixes the marginals once the rotation-sense hidden variable \(\pm\)
is integrated out. We return to this point in
Section~\ref{sec-two-party}; it is the difference between a
stochastic-local model with the right correlator shape and the full
quantum prediction with the right correlator amplitude.

\subsection{The double-angle identity as algebraic
engine}\label{the-double-angle-identity-as-algebraic-engine}

The angle-doubling between linear and quadratic conservation is not a
coincidence of the chosen analyzer geometry. It is the only structural
difference between the two cases.

In the linear case, the conserved quantity
\(\mathbf{p}_A + \mathbf{p}_B\) is \emph{odd} under the symmetry
\(\theta \mapsto \theta + \pi\): the two fragments swap and the sum
changes sign. Under that symmetry, the correlator
\(E(\phi_A, \phi_B) = -\langle \cos(\theta - \phi_A)\cos(\theta - \phi_B)\rangle_\theta\)
inherits a single power of \(\cos\) in each argument, and the integrand
reduces to \(\cos(\Delta\phi)\).

In the quadratic case, the conserved quantity
\(V_\alpha^2 + V_\beta^2 + V_0^2\) is \emph{even} under
\(\theta \mapsto \theta + \pi\): squaring kills the sign change. Under
that symmetry, the correlator inherits a \(\cos^2\) in each argument,
and the double-angle identity converts \(\cos^2(\theta - \phi)\) into
\(\tfrac{1}{2} + \tfrac{1}{2}\cos(2(\theta - \phi))\). The constant term
integrates to a \(\theta\)-independent offset; the
\(\cos(2(\theta - \phi))\) term produces the doubled-frequency
correlator \(\cos(2\Delta\phi)\) after the same orthogonality argument
used in the linear case.

The factor of two is therefore an algebraic invariant of the
conservation law's symmetry class --- specifically, of whether the
conserved quantity is odd or even under emission-phase reversal. It is
not contingent on the detailed dynamics, the analyzer geometry, or the
dimensionality of the underlying field. Any source that emits a
quadratic-conservation-respecting state and is measured by
power-projecting analyzers will reproduce \(\cos(2\Delta\phi)\). Any
source that emits a linear-conservation-respecting state and is measured
by amplitude-projecting analyzers will reproduce \(\cos(\Delta\phi)\).
The Bell--Tsirelson gap is a gap between these two algebraic classes.

\subsection{The hierarchy}\label{the-hierarchy}

Three further models complete the picture. The
\emph{deterministic-threshold model}, in which the continuous quadratic
distribution \(\cos^2(\theta - \phi)\) is converted to a binary outcome
by sign-thresholding, \(A = \mathrm{sgn}[\cos(2(\theta - \phi_A))]\),
produces a piecewise-linear correlator
\(E_\text{det}(\phi_A, \phi_B) = -1 + 4|\Delta\phi|/\pi\) on
\(|\Delta\phi| \leq \pi/2\), with \(S = 2\) exactly --- the Bell bound,
and no violation. The thresholding destroys the cosine structure; the
cos-squared bowl is flattened into a step function and the doubled-angle
information is lost.

The \emph{stochastic-local model}, in which \(\cos^2(\theta - \phi)\) is
interpreted as the probability of a \(+1\) outcome with independent
local coin flips, produces
\(E_\text{stoch}(\phi_A, \phi_B) = -\tfrac{1}{2}\cos(2\Delta\phi)\) and
\(S = \sqrt{2}\). The shape is correct but the amplitude is halved by
the independence of the coin flips: the factorization that defines
locality dilutes the correlation by exactly a factor of two.

The \emph{Popescu--Rohrlich box} (Popescu and Rohrlich 1994), finally,
achieves \(S = 4\) with a discontinuous correlator
\(E(a,b) = -(-1)^{a \cdot b}\) that corresponds to a conservation law so
rigid that the power allocation snaps between channels with zero leakage
to a third mode. No known physical conservation law produces this
behavior, because physical conservation operates on continuous
(Lie-group) symmetries that produce smooth (cosine-type) correlators.

The complete hierarchy is:

\begin{longtable}[]{@{}
  >{\raggedright\arraybackslash}p{(\linewidth - 6\tabcolsep) * \real{0.1959}}
  >{\raggedright\arraybackslash}p{(\linewidth - 6\tabcolsep) * \real{0.2973}}
  >{\raggedright\arraybackslash}p{(\linewidth - 6\tabcolsep) * \real{0.1284}}
  >{\raggedright\arraybackslash}p{(\linewidth - 6\tabcolsep) * \real{0.3784}}@{}}
\toprule\noalign{}
\begin{minipage}[b]{\linewidth}\raggedright
Model
\end{minipage} & \begin{minipage}[b]{\linewidth}\raggedright
\(E(\phi_A, \phi_B)\)
\end{minipage} & \begin{minipage}[b]{\linewidth}\raggedright
\(S_{\max}\)
\end{minipage} & \begin{minipage}[b]{\linewidth}\raggedright
Mechanism
\end{minipage} \\
\midrule\noalign{}
\endhead
\bottomrule\noalign{}
\endlastfoot
Linear conservation & \(-\tfrac{1}{2}\cos(\Delta\phi)\) & \(\sqrt{2}\) &
Single power of \(\cos\); odd-symmetry conservation \\
Stochastic local quadratic & \(-\tfrac{1}{2}\cos(2\Delta\phi)\) &
\(\sqrt{2}\) & Independent coin flips dilute amplitude by
\(\tfrac{1}{2}\) \\
Deterministic threshold & \(-1 + 4\lvert\Delta\phi\rvert/\pi\) & \(2\) &
Binarization linearizes the bowl \\
Quadratic conservation (QM) & \(-\cos(2\Delta\phi)\) & \(2\sqrt{2}\) &
Strict complementarity, \(\cos^2 + \sin^2 = 1\) event-by-event \\
PR box & \(-(-1)^{a \cdot b}\) & \(4\) & Discontinuous (unphysical)
conservation \\
\end{longtable}

Each step in the hierarchy is a factor of \(\sqrt{2}\) in the CHSH
value. The Bell bound \(S = 2\) corresponds to the
deterministic-threshold model, which is the projection of the continuous
quadratic distribution onto binary outcomes with no constraint linking
Alice's and Bob's marginals. The Tsirelson bound \(S = 2\sqrt{2}\)
corresponds to the full quadratic conservation law with strict
event-by-event complementarity --- the constraint that whenever Alice's
analyzer captures power, Bob's analyzer captures the complementary power
on the same trial, enforced not stochastically but rigidly by the
conservation of the quadratic form.

The factor of two between \(S_\text{stoch} = \sqrt{2}\) and
\(S_\text{quad} = 2\sqrt{2}\) is precisely the constraint that the
factorization condition of an LHV model cannot enforce. Independent
local coin flips can reproduce the correct correlator \emph{shape}, but
they cannot enforce the strict complementarity that fixes the correlator
\emph{amplitude}. That fixing requires the global conservation law to
act on every trial. In the present framework, that action is mediated by
the zero-sequence mode --- the third party that is the subject of
Section~\ref{sec-third-party}.

\newpage

\section{The Clarke Transform as Measurement
Apparatus}\label{sec-clarke}

The conservation hierarchy of the previous section gave us the algebraic
shape of the correlator without committing to a specific physical
realization of the source or the analyzer. To go further --- to identify
the third party in the conservation law, to derive the Born-rule
connection, and to construct a classical analog whose statistics
saturate Tsirelson --- we need a concrete representation of ``quadratic
conservation on a vector field with rotational symmetry.'' The Clarke
transform (Clarke 1943) provides exactly this, and its
symmetrical-component decomposition (Fortescue 1918) provides the
natural three-party structure. This section develops the transform, its
power-invariance property, the sequence-component decomposition, and the
Park rotation that implements the analyzer projection. We close with the
geometric correspondence between the Clarke power surface and the
Born-rule probability distribution --- the structural identity that
motivates calling this an analog of single-particle quantum measurement.

\subsection{Definition and
orthonormality}\label{definition-and-orthonormality}

The Clarke transform maps a three-component balanced signal
\((v_a, v_b, v_c)\), with \(v_a + v_b + v_c\) generally nonzero, onto
orthogonal components \((v_\alpha, v_\beta, v_0)\): \[
\begin{pmatrix} v_\alpha \\ v_\beta \\ v_0 \end{pmatrix} \;=\; \sqrt{\frac{2}{3}} \begin{pmatrix} 1 & -\tfrac{1}{2} & -\tfrac{1}{2} \\ 0 & \tfrac{\sqrt{3}}{2} & -\tfrac{\sqrt{3}}{2} \\ \tfrac{1}{\sqrt{2}} & \tfrac{1}{\sqrt{2}} & \tfrac{1}{\sqrt{2}} \end{pmatrix} \begin{pmatrix} v_a \\ v_b \\ v_c \end{pmatrix}.
\] The matrix is orthonormal: its rows are mutually orthogonal unit
vectors, and its inverse is its transpose. The pre-factor \(\sqrt{2/3}\)
is the convention that makes the transform power-invariant (rather than
amplitude-invariant); we adopt it throughout because it is the version
under which Parseval's identity holds without rescaling.

The first two rows define the \(\alpha\) and \(\beta\) axes of a planar
projection of \((v_a, v_b, v_c)\) onto the plane orthogonal to the cube
diagonal \((1,1,1)/\sqrt{3}\). The third row defines the zero-sequence
axis \(v_0\), the projection onto that diagonal. The geometry is:
\(\alpha\beta\) is the plane of ``differential'' mode, \(0\) is the
``common'' mode. A balanced three-phase signal --- one whose three
components sum to zero --- has \(v_0 \equiv 0\) identically. An
imbalanced signal puts energy into the zero-sequence axis. The
decomposition is exhaustive: any three-component signal is uniquely the
sum of a balanced part (in the \(\alpha\beta\) plane) and a common-mode
part (along the \(0\) axis).

\subsection{Power invariance
(Parseval)}\label{power-invariance-parseval}

The defining structural property of the Clarke transform is \emph{power
invariance}. By orthonormality of the transform matrix, \[
v_a^2 + v_b^2 + v_c^2 \;=\; v_\alpha^2 + v_\beta^2 + v_0^2.
\] This is Parseval's theorem applied to the basis change. Total
instantaneous power, computed against a unit resistive load, is
identical in the \(abc\) and \(\alpha\beta 0\) frames. This is the
conservation law that drives everything in the rest of the paper.

The same identity, restated as a sum of three squares equal to a
constant, is the equation of a sphere in three dimensions. The total
available energy budget at any instant is the radius squared. Any
redistribution among the three components --- \(\alpha\), \(\beta\),
\(0\) --- traces out a path on this sphere. The two analyzers in a CHSH
measurement project onto two of the three axes; the third absorbs
whatever the analyzers do not capture. Section~\ref{sec-third-party}
makes this redistribution explicit.

The same conservation law is the statement that quantum-mechanical basis
changes (unitary transformations) preserve \(\langle\psi|\psi\rangle\).
The Clarke transform is the classical, three-dimensional version of that
statement, made on a real vector space rather than a complex Hilbert
space, and acting on power densities rather than probability densities.
This is not analogy. It is the same algebraic identity --- a
sum-of-squares conservation law under an orthonormal change of basis ---
in two different physical realizations.

\subsection{Symmetrical components}\label{symmetrical-components}

The Clarke decomposition, while orthonormal, does not by itself separate
the rotational structure of the signal. A balanced signal in the
\(\alpha\beta\) plane can be written equivalently as a sum of two
counter-rotating phasors --- the \emph{symmetrical components} of
Fortescue (Fortescue 1918). The transformation is \[
\begin{pmatrix} v_+ \\ v_- \\ v_0 \end{pmatrix} \;=\; \frac{1}{3}\begin{pmatrix} 1 & a & a^2 \\ 1 & a^2 & a \\ 1 & 1 & 1 \end{pmatrix} \begin{pmatrix} v_a \\ v_b \\ v_c \end{pmatrix},
\] where \(a = e^{i 2\pi/3}\) is the cube root of unity. The three
sequence components are:

\begin{itemize}
\tightlist
\item
  \textbf{Positive sequence} \(v_+\): phasors rotating
  \emph{counterclockwise} at the source angular frequency. In the
  \(\alpha\beta\) plane, the positive-sequence component traces a circle
  in the \((+)\) direction.
\item
  \textbf{Negative sequence} \(v_-\): phasors rotating \emph{clockwise}.
  In the \(\alpha\beta\) plane, traces a circle in the \((-)\)
  direction.
\item
  \textbf{Zero sequence} \(v_0\): common-mode component. No rotation;
  equal value on all three phases. This is the same \(v_0\) as in the
  Clarke decomposition.
\end{itemize}

The Clarke and Fortescue representations are related by a unitary
transformation on the \(\alpha\beta\) plane: the symmetrical components
are the eigenvectors of the cyclic permutation operator on
\((v_a, v_b, v_c)\), while the Clarke components are the real and
imaginary parts of the positive-sequence phasor (plus the zero-sequence
common mode). The total power, summed across all three sequence
components, equals the power summed across all three Clarke components,
equals the power summed across the original three phases: \[
|v_+|^2 + |v_-|^2 + |v_0|^2 \;=\; v_\alpha^2 + v_\beta^2 + v_0^2 \;=\; v_a^2 + v_b^2 + v_c^2.
\] For a balanced source emitting equal positive and negative sequences
--- the configuration we will use throughout the rest of the paper ---
the zero-sequence component is \emph{initially} zero. It becomes
nonzero, in a sense made precise in Section~\ref{sec-third-party}, only
after the analyzers project the signal onto two specific axes and the
unprojected residue is forced into the common-mode rail by the
conservation constraint.

\subsection{Park rotation as analyzer}\label{park-rotation-as-analyzer}

The analyzer at angle \(\phi\) is implemented as a rotation of the
Clarke frame by \(\phi\) in the \(\alpha\beta\) plane. This is the
\emph{Park transform}: \[
\begin{pmatrix} v_d \\ v_q \end{pmatrix} \;=\; \begin{pmatrix} \cos\phi & \sin\phi \\ -\sin\phi & \cos\phi \end{pmatrix} \begin{pmatrix} v_\alpha \\ v_\beta \end{pmatrix}.
\] The Park rotation does not couple to the zero-sequence axis; it acts
only in the \(\alpha\beta\) plane. The resulting \(d\)-axis (``direct'')
and \(q\)-axis (``quadrature'') components are the projections of the
signal onto the rotated analyzer basis.

For a positive-sequence signal of amplitude \(V\) and emission phase
\(\theta\), the Clarke components are \(v_\alpha = V\cos\theta\) and
\(v_\beta = V\sin\theta\). After the Park rotation by analyzer angle
\(\phi\), the \(d\)-axis power is \[
P_d(\theta, \phi) \;=\; v_d^2 \;=\; \big(v_\alpha\cos\phi + v_\beta\sin\phi\big)^2 \;=\; V^2\cos^2(\theta - \phi).
\] This is Malus's law --- the intensity of polarized light transmitted
through a polarizer at angle \(\phi\). Equivalently, it is the squared
projection of a unit vector at angle \(\theta\) onto an axis at angle
\(\phi\) in the plane. The \(q\)-axis power is the complementary
\(V^2\sin^2(\theta - \phi)\), and the two add to \(V^2\), the total
positive-sequence power available.

The Clarke--Park composition is therefore a \emph{classical polarization
measurement}. The hidden variable \(\theta\) is the emission phase of
the source, distributed uniformly. The analyzer setting \(\phi\) is the
rotation angle of the measurement basis. The output is the squared
projection. This is the same operational structure as a Stern--Gerlach
analyzer on a spin-1/2 particle, a polarizer on a single photon, or a
Mach--Zehnder interferometer on a coherent state, with one important
distinction: every quantity in the Clarke--Park setup is real,
classical, and directly measurable on the signal cables.

\subsection{The Clarke surface and the Born-rule
correspondence}\label{the-clarke-surface-and-the-born-rule-correspondence}

The structural correspondence between the Clarke power decomposition and
quantum single-particle measurement is most visible geometrically.
Following the construction of {[}the present author's notebook 20{]},
consider a balanced source emitting equal positive and negative
sequences with \(V_+ = V_- = V\) and no initial zero-sequence component.
The phase voltages are \[
v_a(\theta) = 2V\cos\theta, \qquad v_b(\theta) = -V\cos\theta, \qquad v_c(\theta) = -V\cos\theta,
\] where the factor of 2 and the sign pattern follow from the
\(120^\circ\) offset cancellation
\(\cos(\theta - 120^\circ) + \cos(\theta + 120^\circ) = -\cos\theta\).
All three phases collapse onto a single sinusoidal envelope: the system
is \emph{linearly polarized} in the \(abc\) frame, its instantaneous
direction set by the relative phase of the two counter-rotating
sequences.

Applying the Clarke transform yields
\(v_\alpha = \sqrt{6}\,V\cos\theta\), \(v_\beta = 0\) identically. The
\(\beta\)-axis is dead. The instantaneous power density in the
\(\alpha\beta\) plane is therefore \[
p_{\alpha\beta}(\theta) \;=\; v_\alpha^2 + v_\beta^2 \;=\; 6V^2\cos^2\theta.
\] Plotted as a radial extrusion, this is a bowl-shaped surface --- the
\emph{Clarke power surface} \(\mathcal{S}\) --- with two peaks at
\(\theta = 0, \pi\) and two nulls at \(\theta = \pm \pi/2\). Three
structural facts about \(\mathcal{S}\) identify it with the Born-rule
probability distribution for a single quantum particle.

\emph{Counter-rotating phasors as particle degrees of freedom.} The
positive sequence rotates counterclockwise; the negative sequence
rotates clockwise. Their interference under the Clarke projection is
maximally constructive along one axis and maximally destructive along
the orthogonal axis, producing the \(\cos^2\theta\) angular
distribution. This is structurally identical to Malus's law for
polarized light, which is itself the Born rule for a single photon's
polarization measurement.

\emph{Power density as Born-rule probability.} The quadratic form
\(p_{\alpha\beta} = v_\alpha^2 + v_\beta^2\), normalized to unit total
energy, is the same object as the Born rule
\(P = |\langle\psi|\phi\rangle|^2\) acting on a two-component complex
amplitude. The amplitude-squared structure is not analogical: it is the
same quadratic form applied to the same kind of object --- a projection
of a two-dimensional state onto an orthonormal basis --- in two
different representations.

\emph{Maximal correlation as axis alignment.} The condition
\(v_\beta \equiv 0\) is the signature of maximal correlation with the
analyzer axis: all power is channeled into one measurement channel, and
none into its orthogonal complement. This is the analog of measuring one
half of a singlet pair along its preparation axis, yielding
\(\cos^2(0) = 1\) in one channel and \(\sin^2(0) = 0\) in the other.

The Clarke surface is the \emph{single-particle marginal} of the joint
Bell object. Tilting the analyzer by \(\phi\) produces the Park rotation
of the previous subsection, which redistributes the power between the
\(d\)-axis bowl (\(\cos^2\phi\) scaling) and the complementary
\(q\)-axis bowl (\(\sin^2\phi\) scaling), with the sum preserved at
\(6V^2\cos^2\theta\) for all \(\phi\). The two bowls \emph{breathe} as
\(\phi\) varies, exchanging amplitude while always summing to the same
constant. The flat plane \(p_d + p_q = 6V^2\cos^2\theta\) is the
conservation law made visible.

The full two-station correlation surface --- the joint object whose CHSH
integral yields \(S = 2\sqrt{2}\) --- is the product of two such
single-particle Clarke surfaces, one per measurement station, with the
joint hidden variable \(\theta\) shared between them. The Bell violation
is then a statement about the integral of this joint object against the
four CHSH analyzer-pair configurations. The next section carries out
that integration explicitly and confronts the threshold/binarization
problem that real Bell tests must solve.

\newpage

\section{Two-Party CHSH from a Three-Phase Source}\label{sec-two-party}

We now derive the CHSH correlation function for a balanced three-phase
source measured by two Park-rotation analyzers, confront the
binarization problem that arises when continuous powers must be
converted to binary \(\pm 1\) outcomes, and identify exactly where each
model in the conservation hierarchy of Section~\ref{sec-hierarchy} sits
in this two-party picture. The central observation is that the two-party
accounting is \emph{complete with respect to local statistics}
(no-signaling holds, marginals are uniform) but \emph{incomplete with
respect to joint statistics} (the conservation law leaks into a third
channel that the two analyzers cannot capture). The leak is the subject
of Section~\ref{sec-third-party}.

\subsection{The standard setup}\label{the-standard-setup}

The source emits equal-amplitude positive and negative sequences,
\(V_+ = V_- = V\), with hidden emission phase \(\theta\) distributed
uniformly on \([0, 2\pi)\) and a hidden rotation-sense label
\(\sigma \in \{+, -\}\) (which sequence the analyzer ``sees'' on a given
trial), distributed uniformly. In the \(\alpha\beta\) frame, the
positive sequence traces a counterclockwise circle of radius \(V\) and
the negative sequence traces a clockwise circle of the same radius.
Their superposition at any instant is a linearly polarized oscillation
along an axis whose direction is set by the relative phase of the two
sequences.

Analyzer \(A\) at angle \(\phi_A\) implements a Park rotation. On a
positive-sequence trial, the captured power is \[
P_A^{(+)}(\theta, \phi_A) \;=\; V^2 \cos^2(\theta - \phi_A).
\] On a negative-sequence trial, the analyzer sees the opposite rotation
sense, and the captured power is the complementary \[
P_A^{(-)}(\theta, \phi_A) \;=\; V^2 \sin^2(\theta - \phi_A).
\] Analyzer \(B\) at angle \(\phi_B\) likewise captures
\(P_B^{(+)} = V^2\cos^2(\theta - \phi_B)\) on positive-sequence trials
and \(P_B^{(-)} = V^2\sin^2(\theta - \phi_B)\) on negative-sequence
trials. The trial-by-trial assignment of rotation sense is the
\emph{anti-correlated} pairing required for the singlet-like state: when
Alice's analyzer sees the positive sequence, Bob's sees the negative,
and vice versa. This is the three-phase analog of the spin assignment in
a singlet decay.

The continuous-valued analyzer outputs, normalized and centered on the
\(\pm 1\) scale used by CHSH, are \[
A(\theta, \phi_A) \;=\; \frac{2 P_A(\theta, \phi_A)}{V^2} - 1 \;=\; \begin{cases} \cos(2(\theta - \phi_A)) & \text{on positive-sequence trials,} \\ -\cos(2(\theta - \phi_A)) & \text{on negative-sequence trials,} \end{cases}
\] and similarly for \(B\), with the rotation-sense pairing enforcing
\(B = -A\) on each trial type.

\subsection{Local uniformity and
no-signaling}\label{local-uniformity-and-no-signaling}

A critical structural feature: neither observer, examining only their
own measurement record, can determine which rotation sense was emitted
on a given trial. The expectation of the analyzer output, taken over the
rotation-sense hidden variable \(\sigma\) at fixed \(\theta\), is \[
\langle A(\theta, \phi_A) \rangle_\sigma \;=\; \tfrac{1}{2}\cos(2(\theta - \phi_A)) + \tfrac{1}{2}\big(\!-\!\cos(2(\theta - \phi_A))\big) \;=\; 0.
\] Integrating over \(\theta\) as well gives \(\langle A \rangle = 0\),
independent of \(\phi_A\). The analyzer's marginal expectation does not
depend on its own setting, and certainly does not depend on the other
analyzer's setting. This is the \emph{no-signaling} condition:
information about Bob's analyzer choice cannot reach Alice through her
measurement record alone.

In density-matrix language, tracing over the rotation-sense hidden
variable yields the maximally mixed state \(\rho_A = I/2\). Local
uniformity is the statement that the source, measured by Alice alone, is
indistinguishable from white noise: every analyzer setting produces the
same uniform marginal distribution. This is the classical statistical
analog of the well-known quantum result that the reduced state of one
half of a Bell pair is the maximally mixed state.

The point bears emphasis. Our system reproduces the Bell correlator at
the joint level (next subsection), but its \emph{marginals are perfectly
classical}. Anything Alice measures, in isolation, looks like a fair
coin. The correlation appears only when her record is compared, after
the fact, with Bob's. This is the same operational signature as quantum
singlets: locally trivial, jointly correlated.

\subsection{The joint correlation}\label{the-joint-correlation}

The joint expectation, taken over the shared hidden variables
\((\theta, \sigma)\), is \[
E(\phi_A, \phi_B) \;=\; \big\langle A(\theta, \phi_A)\, B(\theta, \phi_B) \big\rangle_{\theta, \sigma}.
\] Because the rotation-sense pairing forces \(B = -A\) trial-by-trial,
\[
A(\theta, \phi_A)\, B(\theta, \phi_B) \;=\; -\cos(2(\theta - \phi_A))\,\cos(2(\theta - \phi_B))
\] on both positive- and negative-sequence trials. The rotation-sense
hidden variable drops out of the \emph{joint} expectation, and we are
left with the integral over \(\theta\) alone: \[
E(\phi_A, \phi_B) \;=\; -\int_0^{2\pi} \cos(2(\theta - \phi_A))\,\cos(2(\theta - \phi_B))\, \frac{d\theta}{2\pi} \;=\; -\tfrac{1}{2}\cos(2\Delta\phi),
\] by the same product-to-sum and orthogonality argument used in
Section~\ref{sec-hierarchy}. The factor of \(\tfrac{1}{2}\) is the
amplitude of the correlator before the strict-complementarity constraint
is imposed.

This is the \emph{full continuous-valued correlator} of the three-phase
source. Note that it is the same shape as the quantum singlet correlator
\(E_\text{QM}(\phi_A, \phi_B) = -\cos(2\Delta\phi)\), but with half the
amplitude. The factor of two in the cosine argument is fixed by the
quadratic conservation law; the amplitude is fixed by the binarization
choice, to which we now turn.

\subsection{The binarization problem}\label{the-binarization-problem}

Real CHSH experiments yield binary \(\pm 1\) click outcomes, not
continuous power values. There are two natural ways to binarize a
continuous quadratic distribution, and they correspond to the
\emph{deterministic threshold} and \emph{stochastic local} models of
Section~\ref{sec-hierarchy}. Each gives a CHSH value below the Tsirelson
bound, in the precise way the conservation hierarchy table predicts.

\emph{Deterministic threshold.} Assign \(A = +1\) if
\(\cos^2(\theta - \phi_A) > 1/2\), else \(A = -1\). Equivalently,
\(A(\theta, \phi_A) = \mathrm{sgn}\big[\cos(2(\theta - \phi_A))\big]\),
and similarly
\(B(\theta, \phi_B) = -\mathrm{sgn}\big[\cos(2(\theta - \phi_B))\big]\)
on the anti-correlated rotation-sense pairing. The joint expectation,
integrated over uniform \(\theta\), is the \emph{triangle function} \[
E_\text{det}(\phi_A, \phi_B) \;=\; -1 + \frac{4|\Delta\phi|}{\pi}, \qquad |\Delta\phi| \leq \pi/2,
\] extended periodically. CHSH evaluated at the quadratic-optimal angles
\(\phi_A = 0\), \(\phi_{A'} = \pi/4\), \(\phi_B = \pi/8\),
\(\phi_{B'} = 3\pi/8\) yields \(|S| = 2\) exactly --- the Bell bound,
and no violation. The thresholding has destroyed the cosine structure:
the cos-squared bowl, after binarization, becomes a step function, and
the doubled-angle correlation pattern is replaced by a piecewise-linear
triangle. The doubled-frequency information is gone, and so is any
capacity to violate Bell.

\emph{Stochastic local.} Interpret \(\cos^2(\theta - \phi_A)\) as the
probability of outcome \(A = +1\), conditioned on the hidden variable,
with independent local coin flips at the two stations: \[
\Pr(A = +1 \mid \theta) \;=\; \cos^2(\theta - \phi_A), \qquad \Pr(B = +1 \mid \theta) \;=\; \sin^2(\theta - \phi_B).
\] The locality (factorization) condition is encoded in the independence
of the two coin flips at fixed \(\theta\): \[
\langle AB \mid \theta \rangle \;=\; \langle A \mid \theta\rangle \cdot \langle B \mid \theta \rangle \;=\; \cos(2(\theta - \phi_A)) \cdot \big(\!-\!\cos(2(\theta - \phi_B))\big).
\] Integrating over \(\theta\), \[
E_\text{stoch}(\phi_A, \phi_B) \;=\; -\tfrac{1}{2}\cos(2\Delta\phi), \qquad |S|_\text{stoch} \;=\; \sqrt{2}.
\] The shape is correct, the amplitude is half what the quantum result
would give. The independent coin flips have introduced exactly the
dilution that makes the model \emph{local} in Bell's sense, and the
correlation amplitude has been halved as a consequence.

\subsection{The gap}\label{the-gap}

Three CHSH values are now in hand for the same source statistics,
distinguished only by the binarization rule:

\begin{itemize}
\tightlist
\item
  \emph{Deterministic threshold} \(S = 2\) (Bell bound, no violation;
  piecewise-linear triangle correlator);
\item
  \emph{Stochastic local} \(S = \sqrt{2}\) (correct shape, halved
  amplitude; factorizable);
\item
  \emph{Quantum / strict complementarity} \(S = 2\sqrt{2}\) (Tsirelson
  bound, full amplitude; non-factorizable).
\end{itemize}

Each step in this sequence is a factor of \(\sqrt{2}\) in the CHSH
value. The progression is not arbitrary. The deterministic-threshold
model loses the doubled-frequency structure entirely. The
stochastic-local model preserves the doubled-frequency structure but
halves the correlation amplitude through factorization. The quantum case
preserves both --- and what it requires to do so is the strict
event-by-event complementarity that \[
\cos^2(\theta - \phi_A) + \sin^2(\theta - \phi_A) \;=\; 1 \qquad \text{for every $\theta$ and every $\phi_A$, on every trial.}
\] In the stochastic-local model, this identity holds \emph{in
expectation}: the average power split between Alice's \(d\)-axis and
\(q\)-axis is unity, but on any given trial the binary outcomes need not
respect it. In the quantum case, the same identity holds
\emph{event-by-event}: on every individual trial, the powers in
complementary channels must sum exactly to unity, with no exception.
That is the constraint that the factorization condition cannot enforce,
because factorization requires Alice's outcome to depend only on
\((\phi_A, \theta)\) and Bob's outcome to depend only on
\((\phi_B, \theta)\), and there is no way to make their sum exactly
\(V^2\) on every trial without one outcome conditioning on the other ---
which is precisely the dependence the no-signaling condition forbids on
the marginal level but does not forbid on the joint level.

The doubled-frequency structure --- the cos-squared bowl, the
\(\cos(2\Delta\phi)\) correlator --- is delivered by the quadratic
conservation law on the \emph{source}. The strict-complementarity
amplitude --- the factor of two between \(\sqrt{2}\) and \(2\sqrt{2}\)
--- is delivered by the same quadratic conservation law applied
trial-by-trial across the \emph{full system}, including the channel
between the two analyzers. The two-party accounting captures the first
part of the conservation law (the source emits cos-squared
distributions) but not the second (the joint trials are constrained
event-by-event by a global sum). The missing constraint is the missing
third party, and the next section makes the constraint physical.

\newpage

\section{The Third Party: Zero-Sequence as Vacuum
Mode}\label{sec-third-party}

We now identify the third party in the conservation law. The argument
has three steps. First, we show that the sum of the powers captured by
the two analyzers, \(P_A + P_B\), is \emph{not} constant in \(\theta\);
it oscillates. Second, we show that the missing energy is exactly the
zero-sequence power \(P_0\), which absorbs the residue and restores the
conservation law. Third, we develop the structural analogies between the
zero-sequence mode and (i) the boundary or edge modes of algebraic
quantum field theory, (ii) the quantum potential of Bohmian mechanics,
and (iii) the monogamy of entanglement constraint of Coffman, Kundu, and
Wootters (Coffman, Kundu, and Wootters 2000). The thread connecting the
three is the same: a quadratic conservation law on a bipartite system
has an irreducibly tripartite accounting once the two parts are coupled
through a global constraint, and the third party is the channel through
which the coupling does its work.

\subsection{The missing energy}\label{the-missing-energy}

Take the standard setup of Section~\ref{sec-two-party}: balanced source
emitting equal positive and negative sequences with
\(V_+ = V_- = V = 1\), normalized so total positive-sequence power is
unity. On any trial, with hidden phase \(\theta\) and analyzers at
\(\phi_A\) and \(\phi_B\), the powers captured by the two analyzers in
the positive-sequence channel are \[
P_A(\theta) \;=\; \cos^2(\theta - \phi_A), \qquad P_B(\theta) \;=\; \cos^2(\theta - \phi_B).
\] (For the anti-correlated negative-sequence channel, \(P_B\) is the
complementary \(\sin^2\) form, which we treat in parallel below.) Their
sum, computed using the identity
\(\cos^2 x + \cos^2 y = 1 + \cos(x-y)\cos(x+y)\), is \[
P_A(\theta) + P_B(\theta) \;=\; 1 + \cos(\Delta\phi)\,\cos(2\theta - \phi_A - \phi_B).
\] This is \emph{not} constant in \(\theta\) unless
\(\Delta\phi \in \{0, \pi/2\}\). The two analyzers, taken together, do
not exhaust the energy budget on a generic trial. Some power is captured
by Alice; some is captured by Bob; and some, oscillating in \(\theta\)
at twice the source frequency, is captured by \emph{neither}. That
residue must go somewhere if the total energy is conserved.

In the three-phase analog, where it goes is unambiguous and physically
explicit. The conservation law, restated from Section~\ref{sec-clarke},
is \[
P_+ + P_- + P_0 \;=\; v_a^2 + v_b^2 + v_c^2 \;=\; \mathrm{const},
\] where \(P_+ = |v_+|^2\), \(P_- = |v_-|^2\), and \(P_0 = |v_0|^2\).
The two analyzers project onto the positive-sequence channel (Alice's
\(d\)-axis at \(\phi_A\), Bob's \(d\)-axis at \(\phi_B\)), and what they
fail to capture is forced into the zero-sequence axis by the
conservation constraint. The zero-sequence power, on a generic trial, is
not zero. It is precisely the missing energy: \[
P_0(\theta) \;=\; P_\text{total} - P_A(\theta) - P_B(\theta) \;=\; (P_\text{total} - 1) - \cos(\Delta\phi)\,\cos(2\theta - \phi_A - \phi_B).
\] For a normalized source with \(P_\text{total} = 1\), the constant
offset vanishes and the zero-sequence reservoir oscillates at exactly
the angular frequency required to balance the books.

\subsection{Properties of the zero-sequence
reservoir}\label{properties-of-the-zero-sequence-reservoir}

The behavior of \(P_0(\theta)\) across the four salient analyzer
configurations encodes the entire structure of the CHSH violation.

\emph{Aligned analyzers} (\(\Delta\phi = 0\)). Then
\(\cos(\Delta\phi) = 1\), and \(P_0(\theta) = -\cos(2\theta - 2\phi_A)\)
oscillates with peak-to-peak amplitude two and mean zero. On average,
the two analyzers fully partition the power budget; the
\(\theta\)-by-\(\theta\) oscillation reflects the back-and-forth swing
of the linearly polarized source between full alignment with the
analyzers (everything captured) and full orthogonality (nothing
captured). The captured powers track each other perfectly
trial-by-trial, \(P_A(\theta) = P_B(\theta) = \cos^2(\theta - \phi_A)\),
and the continuous-valued joint correlation
\(E = -\langle (2P_A - 1)(2P_B - 1)\rangle_\theta = -\tfrac{1}{2}\)
attains the maximum magnitude of the half-amplitude correlator.
Strict-complementarity binarization (Section~\ref{sec-two-party})
sharpens this to the singlet result \(E_\text{QM} = -1\) at
\(\Delta\phi = 0\).

\emph{Orthogonal analyzers} (\(\Delta\phi = \pi/2\)). Then
\(\cos(\Delta\phi) = 0\), and \(P_0(\theta) = 0\) identically. The
analyzers' captured powers \(P_A + P_B = 1\) exactly, on every trial,
with no leakage. But the captured powers themselves are also
independent: \(E = 0\), no correlation. The zero-sequence reservoir is
empty, and the two analyzer records are statistically independent. This
is the limit in which the joint state is ``fully partitioned'' between
the two measurement channels with no residue, but at the cost of no
correlation between Alice and Bob.

\emph{CHSH-optimal angles} (\(\Delta\phi = \pi/4, 3\pi/8\), etc.). Then
\(\cos(\Delta\phi)\) is intermediate, and \(P_0(\theta)\) oscillates
with intermediate amplitude. The analyzers capture some of the energy,
the zero-sequence reservoir absorbs the rest. The correlation
\(E = -\tfrac{1}{2}\cos(2\Delta\phi)\) is also intermediate. \emph{This
is precisely the regime in which CHSH violation peaks.} The
analyzer-difference angle that maximizes \(|S|\) is the angle at which
the partition of energy between the captured channels and the
zero-sequence reservoir is in the right balance to produce the strongest
joint correlation.

The CHSH violation, on this reading, is a measurement of the
zero-sequence reservoir. Aligned analyzers drain the reservoir on
average but produce maximum captured-channel correlation. Orthogonal
analyzers empty the reservoir entirely but produce no correlation.
Intermediate angles fill the reservoir to a calibrated level and produce
the maximum CHSH violation. The ``spookiness'' of the Bell violation is,
on this reading, the spookiness of an oscillating power reservoir whose
books are kept by the conservation law but whose channel is invisible to
the standard two-party CHSH apparatus.

\subsection{The three-party conservation
law}\label{the-three-party-conservation-law}

The complete bookkeeping is \[
P_+ + P_- + P_0 \;=\; v_a^2 + v_b^2 + v_c^2 \;=\; v_\alpha^2 + v_\beta^2 + v_0^2 \;=\; \mathrm{const}.
\] This is a \emph{three-party} conservation law. The two analyzers in a
CHSH measurement instrument the first two terms (\(P_+\) and \(P_-\),
accessed via the Park rotation onto the positive- and negative-sequence
channels). The third term (\(P_0\), the zero-sequence reservoir) is
absent from the standard CHSH apparatus.

Standard CHSH analysis, in this language, projects out the third term.
It tracks the two captured channels and treats the residue as
inaccessible noise --- or, more precisely, as the ``spooky'' part of the
joint state that survives only as a post-hoc correlation in the records.
The interpretive move proposed in this paper is to put the third term
back into the bookkeeping. The correlations that appear nonlocal in the
two-party projection become the natural consequence of a global
conservation law on a tripartite system in which the third party is the
channel coupling the two measurement stations.

This is structurally similar to the way momentum conservation in a
two-body decay is \emph{trivially} nonlocal: when one fragment's
momentum is measured, the other's is determined by the conservation law,
regardless of the spatial separation between the measurements. The
strangeness of Bell violations, on the present reading, is not a
categorically different kind of nonlocality. It is the same kind of
nonlocality, applied to a \emph{quadratic} conserved quantity rather
than a \emph{linear} one, and projected onto a two-party measurement
subspace that happens to omit the channel doing the conservation work.
The quadratic structure produces the doubled-frequency correlator (the
\(\cos^2\) bowl) and the projected accounting produces the apparent
surplus correlation (\(S > 2\)) that has no two-party local-realist
explanation.

\subsection{Analogy to algebraic quantum field
theory}\label{analogy-to-algebraic-quantum-field-theory}

The structural role of the zero-sequence mode has a direct analog in
algebraic quantum field theory. In the algebraic formulation (Haag
1996), entanglement between two spatial regions \(A\) and \(B\) is not a
property of particles localized in those regions --- there are, in
general, no such particles in a relativistic field theory --- but a
property of the \emph{state} of the field in a region of spacetime that
contains both. The entanglement entropy between \(A\) and \(B\) is
computed from the modes of the field at the \emph{boundary} between
them, the so-called ``edge modes'' or ``zero modes'' that belong to
neither region's local algebra of observables but that are essential to
the joint state's structure (Witten 2018).

The boundary modes play exactly the role our zero-sequence component
plays. They are the field degrees of freedom that the local observers in
regions \(A\) and \(B\) cannot access through their local algebras, but
whose presence is required to make the global state's bookkeeping
consistent. In the AQFT setting, these modes are responsible for the
divergence of entanglement entropy between adjacent regions and for the
universal area-law behavior at the boundary. In the three-phase analog,
the analogous statement is that the zero-sequence reservoir is the
channel through which the conservation law links Alice's and Bob's
measurement statistics, and its absence from the standard two-party CHSH
bookkeeping is the source of the apparent nonlocality.

The analog is structural rather than dynamical: we are not claiming that
the three-phase system \emph{is} a quantum field theory, only that the
same algebraic pattern --- a quadratic conservation law on a bipartite
measurement system, with a third channel mediating the joint constraint
--- shows up in both, and that the role played by the zero-sequence mode
in our analog is the same role played by edge modes in QFT entanglement.

\subsection{Analogy to Bohmian
mechanics}\label{analogy-to-bohmian-mechanics}

The same structural role is played, in pilot-wave theory (Bohm 1952;
Holland 1993), by the \emph{quantum potential} --- the nonlocal term in
the guidance equation that encodes the global wavefunction's constraint
on local trajectories. In Bohmian mechanics, particles have definite
positions at all times, and their dynamics are determined by a guidance
equation that depends on the wavefunction. The wavefunction acts as a
global object that conditions the local motion: when one particle's
position changes, the wavefunction's structure changes, and the guidance
equation for the other particle is affected even at spacelike
separation. The quantum potential is the explicit nonlocal term that
makes this coupling concrete.

The zero-sequence mode plays the same role. It is the global object
whose state conditions the local powers captured at the two analyzer
stations. When Alice's analyzer setting changes, the partition of energy
between the captured channels and the zero-sequence reservoir changes,
and the joint statistics with Bob's measurements change correspondingly.
The conservation law \(P_+ + P_- + P_0 = \mathrm{const}\) is the
explicit version of the constraint that the quantum potential enforces
implicitly.

This identification suggests a translation table:

\begin{longtable}[]{@{}
  >{\raggedright\arraybackslash}p{(\linewidth - 2\tabcolsep) * \real{0.5000}}
  >{\raggedright\arraybackslash}p{(\linewidth - 2\tabcolsep) * \real{0.5000}}@{}}
\toprule\noalign{}
\begin{minipage}[b]{\linewidth}\raggedright
Bohmian object
\end{minipage} & \begin{minipage}[b]{\linewidth}\raggedright
Three-phase analog
\end{minipage} \\
\midrule\noalign{}
\endhead
\bottomrule\noalign{}
\endlastfoot
Wavefunction & Joint state of the three-mode system
\((v_+, v_-, v_0)\) \\
Particle position & Captured power at an analyzer (\(P_A\) or
\(P_B\)) \\
Quantum potential & Zero-sequence power \(P_0\) \\
Configuration-space nonlocality & Common-mode rail coupling the two
stations \\
Guidance equation & Conservation law with the analyzer rotations as
inputs \\
\end{longtable}

We do not claim that the three-phase system \emph{is} a Bohmian model in
any strong sense --- the dynamics differ, and our framework does not
commit to definite values for unmeasured observables in the way Bohmian
mechanics does. But the structural pattern of ``global constraint
enforcing local correlations through a third channel'' is the same in
both, and the explicit physical realization of the third channel as a
measurable common-mode signal is one of the clarifying features of the
analog.

\subsection{Connection to monogamy of
entanglement}\label{connection-to-monogamy-of-entanglement}

A third structural analogy, less metaphorical and more directly
computable, is to the \emph{monogamy of entanglement} constraint of
Coffman, Kundu, and Wootters (Coffman, Kundu, and Wootters 2000). For a
tripartite quantum system in a pure state, the entanglement of formation
between any pair of subsystems is bounded by a global budget, with the
constraint \[
\mathcal{C}^2(\rho_{AB}) + \mathcal{C}^2(\rho_{AC}) \;\leq\; \mathcal{C}^2(\rho_{A(BC)}),
\] where \(\mathcal{C}\) denotes the concurrence (or analogous
entanglement measure). This is a \emph{budget constraint}: the total
correlation that subsystem \(A\) can share with the rest of the system
is fixed by the global state, and any correlation with \(B\) reduces the
available correlation with \(C\).

The three-party conservation law \(P_+ + P_- + P_0 = \mathrm{const}\) is
the same kind of budget constraint, applied to power rather than
concurrence. The total power available to be distributed across the
three modes is fixed by the source. Any power captured by Alice's
analyzer is power not available to Bob or to the zero-sequence
reservoir. Any power captured by both Alice and Bob is power not
available to the zero-sequence reservoir. The CHSH violation, on this
reading, measures the depth of the trade-off: the analyzers' joint
capture is maximal when the zero-sequence reservoir's depth is
intermediate, and the maximum CHSH value \(S = 2\sqrt{2}\) is the
optimal point on this trade-off curve.

The monogamy constraint and the conservation law are formally analogous,
but the conservation law is more constrained: it requires
\emph{equality}, not inequality, in the budget, on every trial, with the
three modes related by an exact algebraic identity. The monogamy
constraint admits states in which the budget is partially unspent (mixed
states, lossy operations), whereas the pure-state Clarke decomposition
saturates the budget on every trial. Our system is therefore a
\emph{saturating} analog: it sits at the boundary of the monogamy
constraint at all times, which is why it produces Tsirelson saturation
and not some smaller violation.

The three analogies --- AQFT edge modes, Bohmian quantum potential,
monogamy of entanglement --- are not independent. They are three
different views of the same algebraic fact: a quadratic conservation law
on a bipartite measurement system has an irreducibly tripartite
accounting, and the third party is the channel through which the
conservation does its work. In our analog this channel is concrete and
physically measurable; in the abstract quantum frameworks it is implicit
in the wavefunction's structure. The remainder of the paper makes the
three-phase channel explicit at every layer of the construction --- the
reduced 4-mode model in Section~\textbf{?@sec-reduced-model}, the
detector layer in Section~\textbf{?@sec-born}, the closure latch in
Section~\textbf{?@sec-latch}, and the SPICE realization in
Section~\textbf{?@sec-spice} --- and shows that when the channel is
restored to the bookkeeping,

\phantomsection\label{refs}
\begin{CSLReferences}{1}{0}
\bibitem[\citeproctext]{ref-bell1964}
Bell, J. S. 1964. {``On the {Einstein} {Podolsky} {Rosen} Paradox.''}
\emph{Physics Physique Fizika} 1 (3): 195--200.

\bibitem[\citeproctext]{ref-bohm1952}
Bohm, D. 1952. {``A Suggested Interpretation of the Quantum Theory in
Terms of "Hidden" Variables. {I}, {II}.''} \emph{Physical Review} 85:
166--93.

\bibitem[\citeproctext]{ref-clarke1943}
Clarke, E. 1943. \emph{Circuit Analysis of {A-C} Power Systems, Vol.
{I}}. New York: Wiley.

\bibitem[\citeproctext]{ref-chsh1969}
Clauser, J. F., M. A. Horne, A. Shimony, and R. A. Holt. 1969.
{``Proposed Experiment to Test Local Hidden-Variable Theories.''}
\emph{Physical Review Letters} 23 (15): 880--84.

\bibitem[\citeproctext]{ref-coffman2000}
Coffman, V., J. Kundu, and W. K. Wootters. 2000. {``Distributed
Entanglement.''} \emph{Physical Review A} 61 (5): 052306.

\bibitem[\citeproctext]{ref-fortescue1918}
Fortescue, C. L. 1918. {``Method of Symmetrical Co-Ordinates Applied to
the Solution of Polyphase Networks.''} \emph{AIEE Transactions} 37 (2):
1027--1140.

\bibitem[\citeproctext]{ref-giustina2015}
Giustina, M. et al. 2015. {``Significant-Loophole-Free Test of {Bell's}
Theorem with Entangled Photons.''} \emph{Physical Review Letters} 115:
250401.

\bibitem[\citeproctext]{ref-haag1996}
Haag, R. 1996. \emph{Local Quantum Physics: Fields, Particles,
Algebras}. 2nd ed. Springer.

\bibitem[\citeproctext]{ref-hensen2015}
Hensen, B. et al. 2015. {``Loophole-Free {Bell} Inequality Violation
Using Electron Spins Separated by 1.3 Kilometres.''} \emph{Nature} 526:
682--86.

\bibitem[\citeproctext]{ref-hess2001}
Hess, K., and W. Philipp. 2001. {``A Possible Loophole in the Theorem of
{Bell}.''} \emph{Proceedings of the National Academy of Sciences} 98
(25): 14224--27.

\bibitem[\citeproctext]{ref-holland1993}
Holland, P. R. 1993. \emph{The Quantum Theory of Motion: An Account of
the de {Broglie}--{Bohm} Causal Interpretation of Quantum Mechanics}.
Cambridge University Press.

\bibitem[\citeproctext]{ref-thooft2016}
Hooft, G. 't. 2016. \emph{The Cellular Automaton Interpretation of
Quantum Mechanics}. Vol. 185. Fundamental Theories of Physics. Springer.

\bibitem[\citeproctext]{ref-khrennikov2014}
Khrennikov, A. 2014. \emph{Beyond Quantum}. Pan Stanford.

\bibitem[\citeproctext]{ref-larsson2014}
Larsson, J.-Å. 2014. {``Loopholes in {Bell} Inequality Tests of Local
Realism.''} \emph{Journal of Physics A: Mathematical and Theoretical} 47
(42): 424003.

\bibitem[\citeproctext]{ref-popescu1994}
Popescu, S., and D. Rohrlich. 1994. {``Quantum Nonlocality as an
Axiom.''} \emph{Foundations of Physics} 24 (3): 379--85.

\bibitem[\citeproctext]{ref-shalm2015}
Shalm, L. K. et al. 2015. {``Strong Loophole-Free Test of Local
Realism.''} \emph{Physical Review Letters} 115: 250402.

\bibitem[\citeproctext]{ref-spekkens2007}
Spekkens, R. W. 2007. {``In Defense of the Epistemic View of Quantum
States: A Toy Theory.''} \emph{Physical Review A} 75 (3): 032110.

\bibitem[\citeproctext]{ref-tsirelson1980}
Tsirelson, B. S. 1980. {``Quantum Generalizations of {Bell's}
Inequality.''} \emph{Letters in Mathematical Physics} 4 (2): 93--100.

\bibitem[\citeproctext]{ref-witten2018}
Witten, E. 2018. {``{APS} Medal for Exceptional Achievement in Research:
Invited Article on Entanglement Properties of Quantum Field Theory.''}
\emph{Reviews of Modern Physics} 90 (4): 045003.

\end{CSLReferences}




\end{document}
